\subsection{Finite cached data and exact scalar decisions}
\label{sec:newcachebits}

In the static-cache problem, the cache is built before the query is
revealed, and cached coordinates are finite dyadic numbers. This
model already contains the hard instances below.
Section~\ref{sec:models} describes what a cache of transcendental
neural values must provide in addition.

\begin{lemma}[A precision stopping bound]
\label{lem:newprecisionstop}
Let $p\in[0,1]$. Suppose a deterministic routine returns an interval
$I_t\ni p$ of length at most $2^{-t}$ at cost at most
$C(B+t)^d$, for every integer $t\ge1$, where $B\ge1$ and $d\ge1$.
Compare these intervals with the first $t$ bits of one uniform
$U\in[0,1]$, refining until the intervals separate. The resulting
bit is exactly $\mathbf1_{\{U<p\}}$. It stops almost surely and
has expected work at most
\[
 2^{d+5}d!\,C(B+1)^d+O(B+1)
\]
after increasing $C$ by an absolute constant to include interval
arithmetic. The probability of reaching beyond precision $t$ is
at most $\min\{1,2^{2-t}\}$.
\end{lemma}
\paragraph{Proof idea.}
An undecided comparison places the uniform random number close to one fixed probability.  This event shrinks geometrically with precision, and its tail can be summed against every reached arithmetic cost.

\begin{proof}
Replace each interval by its intersection with previous intervals.
The uniform's prefix interval has length $2^{-t}$. If the two
intervals overlap, then $|U-p|\le2^{1-t}$, so the overlap event has
probability at most $2^{2-t}$. Their intersection can persist
forever only when $U=p$, an event of probability zero. Every
decision agrees with the comparison against $p$.

The expected sum of stage costs is at most
\[
 C(B+1)^d+8C\sum_{t\ge2}2^{-t}(B+t)^d.
\]
Since $B\ge1$, $(B+t)^d\le B^d(t+1)^d$. The inequalities
$(t+1)^d\le d!\binom{t+d}{d}$ and the negative-binomial generating
function give
\[
 \sum_{t\ge0}2^{-t}(t+1)^d\le d!2^{d+1}.
\]
The displayed constant is larger than the resulting total. The
random prefix itself uses one additional bit per stage and has
bounded expected length. No derivative under a random stopping
time is used in this argument.
\end{proof}

\begin{lemma}[Exactification from certified intervals]
\label{lem:newexactification}
Suppose a fast query routine supplies an interval for a desired
Bernoulli probability $p$ of width at most $2^{-m}$, at cost $F_m$.
Suppose a fallback routine supplies intervals of width $2^{-t}$
at cost at most $H(B+t)^d$ for every $t>m$.
There is an exact sampler of expected cost
\[
 F_m+O(m)+2^{d+5}d!\,2^{-m}H(B+m+1)^d.
\]
In particular, $m=\max\{1,\lceil\log_2\max\{1,H\}\rceil\}$ makes the
fallback contribution polynomial in $B+\log(H+2)$.
\end{lemma}
\paragraph{Proof idea.}
The fast enclosure settles most uniform comparisons.  A complete evaluation is needed only near its boundary, and choosing the initial precision from the fallback cost makes that rare work affordable.

\begin{proof}
Use the fast interval and $m$ uniform bits first. If the intervals
overlap, continue with the same uniform and the fallback routine
at precisions $m+1,m+2,\ldots$. At stage $m+r$, the probability
of being reached is at most $2^{3-m-r}$. Thus the expected fallback
cost is at most
\[
 8\,2^{-m}H\sum_{r\ge1}2^{-r}(B+m+r)^d.
\]
The same negative-binomial estimate proves the stated bound.
Every interval is certified, so the stopping decision has no error.
\end{proof}

This lemma does not apply merely because an approximation is correct
with high probability. An invalid interval, however rare, can change
the output law. A randomized approximation procedure can be used
when its returned interval is always valid and its expected cost
has the required bound, with that conditional cost accounted for.

\begin{lemma}[A finite-bit scan for cached attention]
\label{lem:newcachescan}
Let $q,k_1,\ldots,k_T$ be dyadic vectors of width $n$, with $b$ bits
per coordinate, and let $v_j\in[-1,1]$ be dyadic. Let $\beta\ge0$
be dyadic of $b$ bits or let $\beta=\sqrt n$.
For $Y=\sum_j e^{a_j}v_j/\sum_j e^{a_j}$, with
$a_j=\beta\langle q,k_j\rangle/n$, the Bernoulli probability
\[
 p=\frac{1+\tanh(\tanh Y)}2
\]
can be sampled exactly with
$O(Tn(b+\log(T+n))^4)$ expected bit operations.
The same asymptotic bound samples an attention index exactly.
\end{lemma}
\paragraph{Proof idea.}
Shift all scores by their maximum, so a denominator term equals one.  Certified absolute errors then control the normalized mean or cumulative probabilities, and interval refinement turns those enclosures into an exact sample.

\begin{proof}
Find a largest inner product using exact dyadic arithmetic.
As $\beta\ge0$, it also maximizes the score. Shift by that score
and put $w_j=e^{a_j-a_{\max}}\le1$. At least one $w_j$ is one,
so $Z=\sum_jw_j\ge1$.

At requested output precision $t$, enclose every $w_j$ to width
$2^{-P}$ with $P=t+\lceil\log_2(T+1)\rceil+8$. The negative
exponential algorithm in Lemma~\ref{lem:newnegativeexp} applies.
Both the numerator and denominator sums then have error at most
$T2^{-P}$. Since the exact numerator has absolute value at most
$Z$ and $Z\ge1$, interval division gives an interval for $Y$ of
width at most $2^{-t-3}$ after clipping to $[-1,1]$.
The two tanh maps are 1-Lipschitz and have bounded arguments.
Their certified series evaluations supply a probability interval
of width at most $2^{-t}$. Its cost is
$O(Tn(b+t+\log(T+n))^4)$, and
Lemma~\ref{lem:newprecisionstop} proves the expected bound.

For an attention index, form intervals for all normalized prefix
weights and compare one lazy uniform with them. There are at most
$T-1$ interior boundaries. Use precision
$P=t+2\lceil\log_2(T+1)\rceil+8$ per weight and read the same
$P$ bits of the uniform. Each boundary enclosure has width at most
$8T2^{-P}$, and the uniform-prefix interval adds at most $2^{-P}$
on either side. The union over the boundaries consequently has
measure at most $10T^2 2^{-P}\le2^{2-t}$. The same stopping
and work calculation applies. This also proves almost-sure
termination if several neighboring weights are extremely small.
\end{proof}

